\VerdierPhysicalPageAnchor{0114}%
\begingroup
\small
\setlength{\abovedisplayskip}{4pt plus 2pt minus 2pt}%
\setlength{\belowdisplayskip}{4pt plus 2pt minus 2pt}%

\VerdierSemanticAnchor{ch2-1-2}
\medskip
\noindent\textbf{1.2. First Properties of Triangulated Categories.}

\VerdierSemanticAnchor{ch2-1-2-1}
\medskip
\noindent\textbf{Proposition 1.2.1.}\quad
Let $\mathcal D$ be a triangulated category. For every object $X$ of
$\mathcal D$, the functors $\operatorname{Hom}_{\mathcal D}(X,\mathord\cdot)$
and $\operatorname{Hom}_{\mathcal D}(\mathord\cdot,X)$ are cohomological
functors with values in the category of abelian groups.

\VerdierSemanticAnchor{ch2-1-2-1-hom-functors}
\noindent\emph{Proof.}\quad
It is enough first to show that the functor
$\operatorname{Hom}_{\mathcal D}(X,\mathord\cdot)$ is cohomological. The
other assertion then follows by passing to the opposite category
$\mathcal D^\circ$ (1.1.7). Thus let

\VerdierSemanticAnchor{ch2-1-2-1-source-triangle}
\[
  X'\xrightarrow{u'}Y'\xrightarrow{v'}Z'\xrightarrow{w'}X'[1]
\]
be a distinguished triangle. It is enough to show that in the diagram

\VerdierSemanticAnchor{ch2-1-2-1-1}
\begin{equation}
  \resizebox{0.98\linewidth}{!}{$
    \operatorname{Hom}_{\mathcal D}(X,X')
    \xrightarrow{\operatorname{Hom}_{\mathcal D}(X,u')}
    \operatorname{Hom}_{\mathcal D}(X,Y')
    \xrightarrow{\operatorname{Hom}_{\mathcal D}(X,v')}
    \operatorname{Hom}_{\mathcal D}(X,Z')
    \xrightarrow{\operatorname{Hom}_{\mathcal D}(X,w')}
    \operatorname{Hom}_{\mathcal D}(X,X'[1])
  $}
  \tag{1.2.1.1}
\end{equation}

\VerdierSemanticAnchor{ch2-1-2-1-exactness-reduction}
we have
\[
  \operatorname{Hom}_{\mathcal D}(X,v')
  \operatorname{Hom}_{\mathcal D}(X,u')=0
\]
and that this diagram is exact at $\operatorname{Hom}_{\mathcal D}(X,Y')$.
Indeed, axiom TRII will then give the same result at
$\operatorname{Hom}_{\mathcal D}(X,Z')$. Thus let
$f\colon X\longrightarrow X'$. By TRI, the triangle

\VerdierSemanticAnchor{ch2-1-2-1-identity-triangle}
\[
  X\xrightarrow{\operatorname{id}_X}X\longrightarrow0\longrightarrow X[1]
\]
is distinguished and, by TRII and TRIII, there exists a morphism
$g\colon X\longrightarrow Y'$ such that the following diagram is commutative:

\VerdierSemanticAnchor{ch2-1-2-1-2}
\begin{equation}
  \begin{tikzcd}[column sep=3.2em,row sep=2.5em]
    X \arrow[r,"{\operatorname{id}_X}"] \arrow[d,"f"']
      & X \arrow[r] \arrow[d,"g"']
      & 0 \arrow[r] \arrow[d]
      & X[1] \arrow[d,"{f[1]}"] \\
    X' \arrow[r,"u'"']
      & Y' \arrow[r,"v'"']
      & Z' \arrow[r,"w'"']
      & X'[1]
  \end{tikzcd}
  \tag{1.2.1.2}
\end{equation}

\VerdierSemanticAnchor{ch2-1-2-1-zero-composite}
It follows that $u'f=g$ and that $v'u'f=v'g=0$, which proves that
$\operatorname{Hom}_{\mathcal D}(X,v')
\operatorname{Hom}_{\mathcal D}(X,u')=0$.

\VerdierSemanticAnchor{ch2-1-2-1-lifting}
Now let $g\colon X\longrightarrow Y'$ be a morphism such that $v'g=0$.
By TRII and TRIII, there exists a morphism $f\colon X\longrightarrow X'$
such that diagram (1.2.1.2) is commutative. It follows that $g=u'f$ and,
consequently, sequence (1.2.1.1) is exact at
$\operatorname{Hom}_{\mathcal D}(X,Y')$.

\VerdierSemanticAnchor{ch2-1-2-2}
\medskip
\noindent\textbf{Corollary 1.2.2.}\quad
In a distinguished triangle
$X'\xrightarrow{u'}Y'\xrightarrow{v'}Z'\xrightarrow{w'}X'[1]$, the
composite of any two consecutive morphisms is zero.

\VerdierSemanticAnchor{ch2-1-2-2-proof}
\noindent\emph{Proof.}\quad
By TRII, it is enough to show that $v'u'=0$. For this, it is enough to take
$X=X'$ in the exact sequence (1.2.1.1) and to consider the successive
images of $\operatorname{id}_{X'}\in\operatorname{Hom}_{\mathcal D}(X',X')$.

\VerdierSemanticAnchor{ch2-1-2-3}
\medskip
\noindent\textbf{Corollary 1.2.3.}\quad
Let $(X,Y,Z,u,v,w)$ and $(X',Y',Z',u',v',w')$ be two distinguished
triangles in $\mathcal D$, and let $(f_1,f_2,f_3)$ be a morphism of
triangles. If two of the $f_i$, $1\leq i\leq3$, are isomorphisms, then so
is the third.
\endgroup

